\section{The one-bit frontier}
\label{sec:onebit}

The smallest interesting record keeps a single bit. There the frontier is exact,
closed-form, and linear in the required margin.

\begin{theorem}[One-bit frontier]
\label{thm:onebit}
Let $n$ be odd, $S$ uniform on $\cube$, $K = 2$, with arbitrary bounded payment
decoders and arbitrary reconstruction. Put
$\Mn = \E[\min(|S|, n-|S|)]$. Then for $0 \le \gamma \le 1/n$,
\[
  \Dopt(2,\gamma) \;=\; \Mn + n\gamma\Bigl(\tfrac{n-1}{2} - \Mn\Bigr),
\]
and no record implements margin $\gamma > 1/n$ at all, whatever the payment.
\end{theorem}

For $K = 2$ the record is a single function $\enc : \cube \to [0,1]$ and
\eqref{eq:margin} factors into a per-agent orientation. Aligning orientations by
input flips---valid precisely because the uniform prior is flip-invariant and
Hamming loss flip-equivariant---turns the feasible set into an affine shift of
the monotone cone, on which the objective is linear for a fixed reconstruction
pair. A linear objective on a shifted cone is what produces a straight frontier.
The optimal record is a mixture of two classical objects, \emph{majority} and a
\emph{random-bit sampler}; neither is new, and the content of the theorem is that
nothing else does better. Appendix~\ref{app:onebit} carries the derivation.

\paragraph{Order of argument (this matters).}
Under \emph{optimal} decoding $D(\enc)$ is a minimum of linear functionals and
hence \emph{concave} in $\enc$, so permutation-averaging can only \emph{increase}
it: symmetrising before the decoder is fixed runs the wrong way for a lower
bound. The correct order, used in Appendix~\ref{app:onebit}, is to fix the
reconstruction pair---for which the objective is linear in $\enc$---optimise over
$\enc$, and only then minimise over pairs. That last step is carried out for
general $n$ by a Rademacher coupling: the term the record can still influence is
$\E[\Delta_{a,b}(S)^-] = \tfrac12\E|R_{p+m}|$, where $R_k$ sums $k$ independent
Rademacher variables and $p+m \le n$ counts the coordinates on which the pair
disagrees. Since $\E|R_k|$ is nondecreasing, every pair is dominated by the
antipodal one.

\begin{remark}[Endpoint]
\label{rem:endpoint}
At $\gamma = 1/n$ the formula collapses to $\Dopt(2,1/n) = (n-1)/2$
independently of $\Mn$: at the maximum implementable margin the record is worth
exactly half a bit per agent less than guessing, whatever $n$ is. The frontier
runs from $(0,\Mn)$ to $(1/n,(n-1)/2)$ and is determined by its endpoints.
\end{remark}

\begin{table}[t]
\centering
\caption{Theorem~\ref{thm:onebit} against an exhaustive-orbit LP, one LP per
reconstruction orbit per $\gamma$ on a grid spanning $[0,1/n]$. Recomputed in the
pinned environment.}
\label{tab:onebit}
\small
\begin{tabular}{@{}lrrr@{}}
\toprule
$n$ & $\Mn$ & $\gamma$ range & max deviation \\
\midrule
3 & 0.750000 & $[0, 1/3]$ & $0$ \\
5 & 1.562500 & $[0, 1/5]$ & $0$ \\
7 & 2.406250 & $[0, 1/7]$ & $4.4\times10^{-16}$ \\
9 & 3.269531 & $[0, 1/9]$ & $3.1\times10^{-15}$ \\
\bottomrule
\end{tabular}
\end{table}

The orbit reduction was cross-checked against complete enumeration of all $4^n$
pairs at $n = 3,5$ (agreement to $0$), and orientation invariance across all
eight sign patterns at $n = 3$ (spread $0$). In every case the optimiser landed
exactly on the conjectured record
$\enc(s) = \gamma|s| + (1-n\gamma)\ind{|s| > n/2}$.
